\section{An infinite-product representation}
\label{a:sec:product}

Take the nonlinear map for the survival probability,
\begin{equation*}
S_{n+1} = 2pS_n - pS_n^2,\qquad S_0=1,
\end{equation*}
and substitute $S_n = 2w_n$. This gives
\begin{equation}
\label{a:eq:logistic}
w_{n+1} = 2p\,w_n(1-w_n) = r\,w_n(1-w_n),
\qquad r = 2p,\quad w_0 = \tfrac12 ,
\end{equation}
which is the logistic map --- the same map that exhibits the period-doubling route
to chaos as $r$ increases from $0$ to $4$ (\cref{a:fig:period_double}). The
subcritical window is $r\in[0,1)$, at the extreme left of that picture, where the
origin is a globally attracting fixed point.

\begin{figure}[ht]
    \centering
    \includegraphics[width=0.86\textwidth]{period_double.png}
    \caption{Bifurcation diagram of the logistic map $w\mapsto rw(1-w)$, showing
    the period doubling of stable states and the onset of chaos. The blue sample
    line on the left is at $r=3.5255$; the right-hand plot shows a trajectory of
    the map at the same $r$.}
    \label{a:fig:period_double}
\end{figure}

In terms of $w$, \cref{a:eq:Apdef} reads
\begin{equation}
\label{a:eq:Aw}
A(p) = \lim_{n\to\infty}\frac{S_n}{(2p)^n} = 2\lim_{n\to\infty}\frac{w_n}{r^n} .
\end{equation}
Any sequence can be written as a telescoping product,
\begin{equation*}
x_n = x_0\lrb{\frac{x_1}{x_0}}\lrb{\frac{x_2}{x_1}}\cdots\lrb{\frac{x_n}{x_{n-1}}}
= x_0\prod_{k=0}^{n-1}\frac{x_{k+1}}{x_k},
\end{equation*}
and applying this with $x_n = w_n/r^n$ the ratio simplifies at once:
\begin{equation}
\label{a:eq:xratio}
\frac{x_{k+1}}{x_k}
= \frac{w_{k+1}}{w_k}\cdot\frac1r
= \frac{rw_k(1-w_k)}{w_k}\cdot\frac1r
= 1-w_k .
\end{equation}
The whole of the nonlinearity has been isolated into a single factor per
generation, and the geometric decay has cancelled out entirely.

\begin{proposition}[Product representation]
\label{a:prop:product}
For $p\in(0,\tfrac12)$, with $S_n$ generated by \cref{a:eq:SnIter} and
$w_k = S_k/2$,
\begin{equation}
\label{a:eq:prodFormula}
A(p) = \prod_{k=0}^{\infty}(1-w_k)
= \frac12\prod_{k=1}^{\infty}\lrb{1-\frac{S_k}{2}} ,
\end{equation}
and the product converges to a strictly positive limit. In particular the limit
\cref{a:eq:Apdef} defining $A(p)$ exists, and $0<A(p)<\tfrac12$.
\end{proposition}

\begin{proof}
Iterating \cref{a:eq:xratio} gives $x_n = x_0\prod_{k=0}^{n-1}(1-w_k)$ with
$x_0 = w_0 = \tfrac12$. For $0<w_k<1$ the infinite product $\prod_{k\ge1}(1-w_k)$
converges to a strictly positive limit if and only if $\sum_{k\ge1}w_k<\infty$.
From \cref{a:eq:logistic} with $0<r<1$ and $0<w_n<1$ we have $w_{n+1}<rw_n$, so by
induction $w_n\le w_0r^n$; since $0<r<1$ this bound is summable, and the product
converges to a positive finite value. Hence $\lim_nx_n$ exists and is positive, and
by \cref{a:eq:Aw},
\begin{equation*}
A(p) = 2\lim_nx_n = 2\cdot\tfrac12\prod_{k=0}^{\infty}(1-w_k)
= \prod_{k=0}^{\infty}(1-w_k) .
\end{equation*}
Pulling out the factor $1-w_0 = \tfrac12$ gives the second form, and since every
remaining factor lies strictly in $(0,1)$ for $p>0$, the value is strictly below
$\tfrac12$.
\end{proof}

\subsection{A certified truncation}
\label{a:sec:producttail}

The partial products
\begin{equation}
\label{a:eq:partialproduct}
A_N(p) = \frac12\prod_{k=1}^{N}\lrb{1-\frac{S_k}{2}}
\end{equation}
decrease monotonically to $A(p)$, so every truncation is a rigorous upper bound. A
matching lower bound follows from the elementary inequality
$\prod_k(1-a_k)\ge1-\sum_ka_k$ for $a_k\in[0,1]$ applied to the discarded tail.
Since $w_{k+1}\le rw_k$, the tail sum is dominated by a geometric series and
\begin{equation}
\label{a:eq:producttailbound}
A_N(p)\lrb{1-\frac{w_{N+1}}{1-r}} \;\le\; A(p) \;\le\; A_N(p),
\end{equation}
whenever the left-hand factor is positive.

\subsection{Is this an advance?}
\label{a:sec:isadvance}

Only a modest one. \Cref{a:eq:prodFormula} is exact, it proves existence, and by
\cref{a:eq:producttailbound} it brackets $A(p)$ at every order. But computing the
factors still requires iterating the
nonlinear map to generate the $S_k$, so the near-critical difficulty is untouched:
as $r\to1$ the bound $w_n\le w_0r^n$ degrades and the number of factors needed
grows without limit.
